% Nouveau EGA V (cote 25): archive cote 25, site page 113 (PDF page 114).
% typescript. Draft compared: none.
% AI-transcribed by Claude Opus 5.5 (claude-opus-5-5, adaptive thinking, effort high), Message Batches API, batch ega25b (msgbatch_01A7n5YuXubUPNK1zsepi8RG), 2026-11-13; not yet checked by eye.
\scanpage{25/113}
\mspage{\hand{16}}

\var{\xout{\unsure{Parxsymétxiex}}}{}{typewriter strike-out}{ega25_25_113-v0001} Lorsque Z est fermé, alors X-Z est le plus grand ouvert
tel que la famille de diviseurs de X-Z paramétrée par T, induite par la
famille donnée (dans un sens évident) est sans point fixe.

Si la famille D est sans point fixe, et si T est \var{\inserted{plat}}{plat}{insertion}{ega25_25_113-v0002} loc de prés finie
sur S à fibres $(S_1)$\var{\inserted{et géom irréd}}{et géom irréd}{insertion}{ega25_25_113-v0003}, alors D est aussi un diviseur \var{\deleted{de Cartier}}{}{deletion}{ega25_25_113-v0004}
relativement à X (pour $pr_1 : X\times_S T \longrightarrow X$): en effet, D est défini localemen
% NOTE: the right edge of the page is cut; the end of this line and of several following lines is not visible
\var{\inserted{au point $z \in$ supp D}}{au point $z \in$ supp D}{insertion}{ega25_25_113-v0005} par une équation \var{\xout{\unsure{f=0}}}{}{typewriter strike-out}{ega25_25_113-v0006} $\phi = 0$, et l'équation induite sur la fibre
$T_{s\ k(x)}$ du point $x\in X$ (au dessus de $s\in S$) est non \var{\xout{\unsure{identiquement nulle}}}{}{typewriter strike-out}{ega25_25_113-v0007}
nilpo\var{\xout{\unsure{n}}}{}{typewriter strike-out}{ega25_25_113-v0008}tente en z (car autrement $pr_1^{-1}(x)\ D = V(\phi_x)$ \var{\xout{\unsure{serait}}}{}{typewriter strike-out}{ega25_25_113-v0009} \var{\inserted[typed]{contiendrait}}{contiendrait}{insertion}{ega25_25_113-v0010} ensemblist\var{\xout{\unsure{e}}}{}{typewriter strike-out}{ega25_25_113-v0011}
ment un voisinage de z dans $T_{\var{\xout{s}}{}{typewriter strike-out}{ega25_25_113-v0012}\ k(x)}$ , donc contiendrait $T_{s\ k(x)}$ , i.e.
x serait point fixe, ce qui n'est pas);or comme $T_{s\ k(x)}$ est irréd et $(S_1)$
il s'ensuit que $\phi_x$ est $\underline{O}_{T_{s\ k(x)}}$-régulière en z . On obtient ainsi
\var{\xout{\unsure{un morphisme}}}{}{typewriter strike-out}{ega25_25_113-v0013} une famille de div de T/S paramétrée par X, i.e. un morph
\[ X \longrightarrow \mathrm{Div}_{T/S} \quad , \]
\hand{D}ans le cas général où la famille de diviseurs de X peut avoir des points
fixes, on obtient une famille de diviseurs de T/S paramétrée par X-Z,
i;e. un morphisme
\[ X - Z \longrightarrow \mathrm{Div}_{T/S} \ , \]
en remplaçant dans la définition précédente X pa\var{\xout{r}}{}{typewriter strike-out}{ega25_25_113-v0014} X-T . D'ailleurs la
\var{\xout{si D est sans points fixes,}}{}{typewriter strike-out}{ega25_25_113-v0015}

\var{\xout{Notons en passant que la famille de diviseurs (notée ${}^{t}D$ )}}{}{typewriter strike-out}{ega25_25_113-v0016}
\var{\xout{X-Z}}{}{typewriter strike-out}{ega25_25_113-v0017}
\var{\xout{de T/S paramétrée par X est}}{}{typewriter strike-out}{ega25_25_113-v0018} \var{\deleted{évidemment sans points fixes, et on}}{}{deletion}{ega25_25_113-v0019}
\var{\xout{obtient ainsi une correspondance biunivo}}{}{typewriter strike-out}{ega25_25_113-v0020}\var{\deleted{que entre familles de diviseurs}}{}{deletion}{ega25_25_113-v0021}
\var{\xout{sans points fixes}}{}{typewriter strike-out}{ega25_25_113-v0022} démonstration ci-dessus montre que X-Z est exactement
le plus grand ouvert \var{\inserted{U}}{U}{insertion}{ega25_25_113-v0023} de X tel que $D|(U\times_S T)$ soit un di\var{\xout{vi}}{}{typewriter strike-out}{ega25_25_113-v0024}seur \var{\deleted{de Cartier}}{}{deletion}{ega25_25_113-v0025}
relativement à U, i.e. \var{\inserted{tel que son symétrique ${}^{t}D$}}{tel que son symétrique ${}^{t}D$}{insertion}{ega25_25_113-v0026} \unsure{``soit''} une famille de diviseurs de T/S paramétrée
par U/S .

\mspage{\hand{113}}
